% !TEX program = xelatex
% ============================================================
%  خطة تنسيق الألوان حسب الثيمات — صفحة الإعدادات
%  (الوضع الفاتح / الداكن / التراثي-سيپيا)
%  الاستخدام على Overleaf:
%    1) New Project ← Blank Project
%    2) الصق المحتوى في main.tex
%    3) Menu ← Compiler ← اختر: XeLaTeX   (مهم للعربية)
%    4) Recompile ← نزّل PDF
% ============================================================
\documentclass[11pt]{article}

\usepackage{polyglossia}
\setmainlanguage{arabic}
\setotherlanguage{english}
\newfontfamily\arabicfont[Script=Arabic,Scale=1.0]{Amiri}
\newfontfamily\arabicfontsf[Script=Arabic,Scale=1.0]{Amiri}
\setmonofont{DejaVu Sans Mono}[Scale=0.8]

\usepackage{geometry}
\geometry{a4paper, margin=2.2cm}

\usepackage{xcolor}
\usepackage{graphicx}
\usepackage{booktabs}
\usepackage{longtable}
\usepackage{array}
\usepackage{enumitem}
\usepackage{titlesec}
\usepackage{fancyhdr}
\usepackage{hyperref}
\usepackage{tcolorbox}
\tcbuselibrary{skins,breakable}
\usepackage{fontawesome5}

% ── الألوان (هوية يَسِّر) ──
\definecolor{gold}{HTML}{B8933B}
\definecolor{teal}{HTML}{0F3D3E}
\definecolor{darkbg}{HTML}{101414}
\definecolor{okgreen}{HTML}{27AE60}
\definecolor{warnred}{HTML}{C0392B}
\definecolor{softgray}{HTML}{F4F4F4}
\definecolor{midgray}{HTML}{7A7A7A}
\definecolor{codebg}{HTML}{1A2020}
\definecolor{sepiapaper}{HTML}{F4ECD8}

\pagestyle{fancy}
\fancyhf{}
\renewcommand{\headrulewidth}{0.4pt}
\renewcommand{\footrulewidth}{0.4pt}
\fancyhead[R]{\small\color{gold}خطة تنسيق الألوان حسب الثيمات — منصة يَسِّر}
\fancyhead[L]{\small\color{midgray}Theme Color System}
\fancyfoot[C]{\small\thepage}

\titleformat{\section}{\Large\bfseries\color{gold}}{\thesection}{0.6em}{}
\titleformat{\subsection}{\large\bfseries\color{darkbg}}{\thesubsection}{0.5em}{}
\titleformat{\subsubsection}{\normalsize\bfseries\color{midgray}}{\thesubsubsection}{0.4em}{}

\setlist[itemize]{rightmargin=0pt, itemsep=2pt, topsep=3pt}
\setlist[enumerate]{rightmargin=0pt, itemsep=2pt, topsep=3pt}

\hypersetup{
  colorlinks=true, linkcolor=gold, urlcolor=gold,
  pdftitle={خطة تنسيق الألوان حسب الثيمات}, pdfauthor={منصة يَسِّر}
}

% ── صناديق مخصّصة ──
\newtcolorbox{goalbox}[1][]{
  colback=gold!8, colframe=gold, boxrule=0.8pt,
  arc=3pt, left=8pt, right=8pt, top=6pt, bottom=6pt,
  breakable, #1
}
\newtcolorbox{notebox}[1][]{
  colback=softgray, colframe=midgray, boxrule=0.6pt,
  arc=3pt, left=8pt, right=8pt, top=6pt, bottom=6pt,
  breakable, #1
}
\newtcolorbox{warnbox}[1][]{
  colback=warnred!6, colframe=warnred, boxrule=0.8pt,
  arc=3pt, left=8pt, right=8pt, top=6pt, bottom=6pt,
  breakable, #1
}
\newtcolorbox{codebox}[1][]{
  colback=codebg, colframe=codebg, coltext=white,
  arc=3pt, left=8pt, right=8pt, top=6pt, bottom=6pt,
  breakable, #1
}

\begin{document}

% ============================================================
\begin{center}
  {\Huge\bfseries\color{gold} خطة تنسيق الألوان حسب الثيمات}\\[4pt]
  {\large\color{darkbg} توحيد صفحة الإعدادات على نظام الثيمات الثلاثة}\\[6pt]
  {\small\color{midgray} منصة يَسِّر الفقهية — \today}
\end{center}

\vspace{4pt}
\hrule
\vspace{10pt}

% ============================================================
\section{الهدف}

\begin{goalbox}
إخضاع صفحة الإعدادات بالكامل لنظام الثيمات الموجود في المنصة
(\textbf{فاتح} / \textbf{داكن} / \textbf{تراثي-سيپيا}) عبر متغيّرات CSS، وإزالة الألوان الثابتة
(hardcoded) التي لا تتغيّر مع تبديل الثيم، مع ضمان تباين لوني كافٍ (WCAG AA) في كل وضع.
\end{goalbox}

% ============================================================
\section{نظام الثيمات الحالي}

الثيمات تُطبَّق عبر \texttt{data-theme} على عنصر \texttt{<html>}، وتعريف المتغيّرات في
\texttt{src/index.css}:

\vspace{6pt}
\renewcommand{\arraystretch}{1.35}
\begin{longtable}{@{}p{0.30\textwidth} p{0.22\textwidth} p{0.22\textwidth} p{0.20\textwidth}@{}}
\toprule
\textbf{المتغيّر} & \textbf{فاتح} & \textbf{داكن} & \textbf{سيپيا} \\
\midrule
\endhead
\texttt{--bg-color}        & \texttt{\#FAFAF7} & \texttt{\#101414} & \texttt{\#F4ECD8} \\
\texttt{--primary-color}   & \texttt{\#0F3D3E} & \texttt{\#5CC0B5} & \texttt{\#5B3A29} \\
\texttt{--accent-color}    & \texttt{\#FBDC99} & \texttt{\#D8A06A} & \texttt{\#C09553} \\
\texttt{--text-main}       & \texttt{\#2C2C2C} & \texttt{\#E6E8E6} & \texttt{\#3E2A1A} \\
\texttt{--text-muted}      & \texttt{\#6C757D} & \texttt{\#A9B1AD} & \texttt{\#7A6A5A} \\
\texttt{--card-bg}         & \texttt{\#FFFFFF} & \texttt{\#1A2020} & \texttt{\#FAF3E3} \\
\texttt{--border-color}    & \texttt{\#EEEEEE} & \texttt{rgba(...0.09)} & \texttt{rgba(...0.15)} \\
\texttt{--badge-bg}        & \texttt{\#F8F9FA} & \texttt{\#232B2B} & \texttt{\#F0E6D0} \\
\texttt{--text-on-accent}  & \texttt{\#0F3D3E} & \texttt{\#121212} & \texttt{\#3E2A1A} \\
\bottomrule
\end{longtable}
\renewcommand{\arraystretch}{1.0}

% ============================================================
\section{المشكلة في الوضع الحالي}

\begin{warnbox}
ملف \texttt{SettingsView.module.css} يحتوي قرابة \textbf{75 لونًا ثابتًا} لا تتبع الثيم، منها:
\begin{itemize}
  \item \textbf{الذهبي/التيل الثابت:} \texttt{rgba(201,168,76,...)} و \texttt{rgba(92,192,181,...)} في
        ترويسة الصفحة وأيقونات البطاقات — لا تتبدّل مع ثيم سيپيا أو الداكن.
  \item \textbf{ألوان الحالات الثابتة:} الأخضر \texttt{\#27AE60}، الأحمر \texttt{\#E74C3C}، الأزرق
        \texttt{\#2980B9}، البرتقالي \texttt{\#D35400} — مضبوطة يدويًا في موضعين لكل حالة (فاتح/داكن)
        دون شمل ثيم سيپيا.
  \item \textbf{الألوان الاجتماعية:} واتساب/فيسبوك/لينكد إن بعلامتها التجارية — مقبولة لكن تحتاج ضبطًا
        للخلفية في الوضع الداكن.
  \item \textbf{الظلال:} \texttt{rgba(0,0,0,...)} ثابتة تكون باهتة/ثقيلة في الوضع الداكن.
\end{itemize}
\end{warnbox}

\begin{notebox}
ثيم \textbf{سيپيا} معرَّف في \texttt{index.css} لكنه \textbf{لا يحصل على أي ضبط خاص} في صفحة الإعدادات،
فتظهر الألوان نفسها المخصّصة للوضع الفاتح، ما يخلّ باتّساق المظهر التراثي.
\end{notebox}

% ============================================================
\section{المرحلة الأولى: إضافة رموز دلالية (Design Tokens)}

إضافة مجموعة متغيّرات جديدة في \texttt{:root} داخل \texttt{index.css} للدلالات اللونية
(نجاح/تحذير/خطأ/معلومة/ذهبي-هوية)، ثم تعريفها في كل ثيم:

\vspace{6pt}
\renewcommand{\arraystretch}{1.35}
\begin{longtable}{@{}p{0.24\textwidth} p{0.24\textwidth} p{0.24\textwidth} p{0.22\textwidth}@{}}
\toprule
\textbf{الرمز الدلالي} & \textbf{فاتح} & \textbf{داكن} & \textbf{سيپيا} \\
\midrule
\endhead
\texttt{--tone-success}  & \texttt{\#27AE60} & \texttt{\#2ECC71} & \texttt{\#3F7A46} \\
\texttt{--tone-warning}  & \texttt{\#D35400} & \texttt{\#F39C12} & \texttt{\#B0762A} \\
\texttt{--tone-danger}   & \texttt{\#E74C3C} & \texttt{\#FF6B5B} & \texttt{\#B23A2E} \\
\texttt{--tone-info}     & \texttt{\#2980B9} & \texttt{\#5DADE2} & \texttt{\#4A6FA5} \\
\texttt{--brand-gold}    & \texttt{\#C9A84C} & \texttt{\#D8A06A} & \texttt{\#C09553} \\
\texttt{--brand-teal}    & \texttt{\#5CC0B5} & \texttt{\#5CC0B5} & \texttt{\#7A9E8E} \\
\texttt{--shadow-soft}   & \texttt{rgba(0,0,0,.06)} & \texttt{rgba(0,0,0,.35)} & \texttt{rgba(90,60,30,.10)} \\
\bottomrule
\end{longtable}
\renewcommand{\arraystretch}{1.0}

\begin{goalbox}
وبهذا تحتاج ملفات المكوّنات إلى استدعاء رمز واحد \texttt{var(--tone-success)} بدل كتابة قيمتين
(فاتح/داكن) لكل حالة، وحلّ ثيم سيپيا تلقائيًا.
\end{goalbox}

% ============================================================
\section{المرحلة الثانية: تحويل الألوان الثابتة}

خريطة الاستبدال في \texttt{SettingsView.module.css}:

\vspace{6pt}
\renewcommand{\arraystretch}{1.35}
\begin{longtable}{@{}p{0.42\textwidth} p{0.46\textwidth}@{}}
\toprule
\textbf{القيمة الحالية الثابتة} & \textbf{البديل الموصى به} \\
\midrule
\endhead
\texttt{rgba(201,168,76,.18)} (ذهبي الهوية) & \texttt{color-mix(in srgb, var(--brand-gold) 18\%, transparent)} \\
\texttt{rgba(92,192,181,.22)} (تيل الهوية) & \texttt{color-mix(in srgb, var(--brand-teal) 22\%, transparent)} \\
\texttt{\#27AE60} / \texttt{\#2ECC71} & \texttt{var(--tone-success)} \\
\texttt{\#D35400} / \texttt{\#F39C12} & \texttt{var(--tone-warning)} \\
\texttt{\#E74C3C} & \texttt{var(--tone-danger)} \\
\texttt{\#2980B9} / \texttt{\#5DADE2} & \texttt{var(--tone-info)} \\
\texttt{rgba(0,0,0,.05..0.08)} (الظلال) & \texttt{var(--shadow-soft)} \\
\texttt{\#C99347} (تدرّج الزر) & \texttt{color-mix(in srgb, var(--accent-color) 85\%, #000)} \\
\bottomrule
\end{longtable}
\renewcommand{\arraystretch}{1.0}

\begin{notebox}
\texttt{color-mix()} مدعوم في كل المتصفحات الحديثة (Chrome/Edge/Safari/Firefox).
كبَديل آمن للتوافق، تُعرّف شفافيات صريحة في كل ثيم:
\texttt{--brand-gold-18: rgba(...)}. اختر أحدهما ووحّده في الملف.
\end{notebox}

% ============================================================
\section{المرحلة الثالثة: ضبط خاص لكل ثيم}

\subsection*{الوضع الداكن}
\begin{itemize}
  \item الظلال تزداد قوة (\texttt{rgba(0,0,0,.35)}) لتعويض الخلفية الداكنة.
  \item الخلفيات الشفافة تُرفع نسبتها (من $0.10$ إلى $0.16$) لتبقى مرئية على التيل الداكن.
  \item ألوان العلامات التجارية (واتساب...) يُخفّف تشبّعها نحو $0.9$ لتفادي الإبهار.
\end{itemize}

\subsection*{ثيم سيپيا (تراثي)}
\begin{itemize}
  \item المشبّع يُقلّل والأشكال تُدفأ: الأخضر يميل للزيتي، والأزرق يميل للكحلي الورقي.
  \item الذهبي يصبح \texttt{\#C09553} (مساويًا \texttt{--accent-color} في سيپيا).
  \item الحدود تُبنى على \texttt{rgba(90,60,30,.15)} بدل الأسود/الأبيض.
\end{itemize}

% ============================================================
\section{المرحلة الرابعة: إتاحة ومعايير}

\begin{itemize}
  \item الالتزام بتباين \textbf{WCAG AA} (نسبة $4.5{:}1$ للنص العادي) في كل ثيم.
  \item عدم استخدام اللون وحده لنقل المعنى: كل حالة تحمل أيقونة ونصًا أيضًا.
  \item اختبار فعلي بأداة تباين على الأزواج: (نص/خلفية)، (شارة/خلفية البطاقة).
  \item احترام \texttt{prefers-color-scheme} عند الثيم الافتراضي إن رُغب مستقبلًا.
\end{itemize}

% ============================================================
\section{الملفات والتحقق}

\textbf{ملفات تُعدَّل:}
\begin{itemize}
  \item \texttt{src/index.css} — إضافة الرموز الدلالية في الثيمات الثلاثة.
  \item \texttt{src/Components/Views/SettingsView.module.css} — استبدال الألوان الثابتة بالمتغيّرات.
\end{itemize}

\textbf{لا يُعدَّل:} بنية JSX ولا منطق المكوّن.

\vspace{6pt}
\textbf{أوامر التحقق:}
\begin{codebox}
{\ttfamily\small\raggedleft
npm run lint\\
npm run test -{}- SettingsView\\
npx tsc -{}-noEmit
}
\end{codebox}

\begin{warnbox}
\textbf{ملاحظة توافق:} عند استخدام \texttt{color-mix()} سيبقى جاهزًا؛ أما إن اخترت النهج البديل
فأضِف الشفافيات الصريحة إلى الثيمات الثلاثة قبل حذف القيم القديمة لتجنّب أي كسر بصري مؤقت.
\end{warnbox}

% ============================================================
\section{معايير القبول}

\begin{itemize}
  \item[${\color{okgreen}\faCheckCircle}$] لا يوجد أي لون ثابت (غير متغيّر) في \texttt{SettingsView.module.css}.
  \item[${\color{okgreen}\faCheckCircle}$] الثيمات الثلاثة تعرض حالة متّسقة ومقروءة لصفحة الإعدادات.
  \item[${\color{okgreen}\faCheckCircle}$] شارات الحالة ورسائل التغذية الراجعة مقروءة في كل ثيم.
  \item[${\color{okgreen}\faCheckCircle}$] ثيم سيپيا يحصل على ضبطه الخاص (لا يرث مظهر الفاتح).
  \item[${\color{okgreen}\faCheckCircle}$] نسبة التباين تستوفي WCAG AA للحالات الأساسية.
\end{itemize}

\end{document}
